Agent-memory benchmarks score a system by how often its agent answers questions about the past correctly. But when a
memory drops a fact, the underlying language model can often still produce the right answer from its prior, so the
score mixes remembering with guessing. We make this measurable. Closed-book, the same reader answers
\gLocomoPct{} of LoCoMo and \gQaOnePct{} of MemoryAgentBench retrieval items correctly in expectation, and up to \everQaOne{}
of the latter at least once in seven draws. We propose \pcme{}, an evaluation procedure that estimates each item's
guess rate with the evaluation prompt itself, reports accuracy on prior-resistant items (\cpr{}) as the primary memory
score, defines a guess-corrected score under a forced-answer protocol, and scores unanswerable items by discrimination
rather than accuracy. Each estimand comes with its identifying assumptions and a pre-registered diagnostic, and we
report where the diagnostics failed. Re-scoring twelve memory systems shows that differences in raw accuracy between
memory-poor systems are driven by their tendency to abstain (Spearman $\rho=\hOneQaOne$ between abstention and the gain
from forcing an answer), and that LoCoMo's unanswerable category rewards forgetting ($\rho=\rhoCaCpr$ with \cpr{}). A
controlled study with \nFinal{} counter-schema minimal pairs, replicated with two model families, shows the mechanism:
memories delete prior-congruent and prior-incongruent facts alike, and the agent then fills the gaps from its prior, so
belief gaps reach $\dKrwQwenHi$ while raw accuracy looks healthy.
